% !PRINTED 122
\weberpage{122}
%<<CONT
hier nach einem Verfahren von \gesperrt{Frobenius} beweisen, das sich
nur auf den allgemeinen Gruppenbegriff stützt\footnote{\gesperrt{Cauchy}, Exerc. d'analyse, Tom. III. \gesperrt{Sylow}, Mathem. Annalen, Bd. 5. \gesperrt{Frobenius}, Journ. für Mathematik, Bd. 100. \gesperrt{Netto}, Mathem. Annalen, Bd. 13; Substitutionentheorie. \S.\,48.}.

Der Satz lässt sich einfach so aussprechen:

\gesperrt{I. Ist $P$ eine Gruppe von irgend welchen Elementen vom Grade $n$ und $p^{\varkappa}$ eine in $n$ aufgehende Primzahlpotenz, so hat die Gruppe $P$ einen Theiler vom Grade $p^{\varkappa}$.}

\gesperrt{Cauchy} hat diesen Satz für $\varkappa = 1$ bewiesen. Für den Fall, dass $n$ eine Primzahl ist, ist er evident; für $n = 4$ und $n = 6$ kann man ihn aus den Zusammenstellungen in \S.\,27 leicht ablesen, so dass er also für die Fälle $n = 2, 3, 4, 5, 6, 7$ als erwiesen betrachtet werden kann. Auf Grund dieser Wahrnehmung lässt sich die vollständige Induction anwenden, und wir setzen also voraus, \gesperrt{der Satz sei bewiesen für Gruppen, deren Grad niedriger ist als} $n$.

Wir wollen die Elemente der Gruppe $P$ durch die Buchstaben $a, b, c, \ldots$ bezeichnen und unter $x$ ein Zeichen für ein veränderliches Element in $P$ verstehen, das die ganze Gruppe $P$ durchläuft.

Wir fassen zunächst alle die Elemente $a$ zusammen, die der Bedingung
\weq{1}{x^{-1}ax = a \quad \text{oder} \quad ax = xa}
genügen, so dass $a$ jedes Element bedeutet, das mit allen Elementen von $P$ \gesperrt{vertauschbar} ist. Dazu gehört jedenfalls das Einheitselement, und die Gesammtheit dieser Elemente $a$ bildet eine in $P$ enthaltene Gruppe $A$. Denn aus
\[ x^{-1}ax = a,\quad x^{-1}a'x = a' \]
folgt
\[ x^{-1}aa'x = aa'. \]

\gesperrt{$A$ ist eine Abel'sche Gruppe und ein Normaltheiler von $P$}; denn aus der Definition (1) folgt $x^{-1}Ax = A$, und für irgend zwei Elemente $a, a'$ von $A$:
\[ aa' = a'a. \]
Wir bezeichnen mit $\nu$ und $j$ Grad und Index dieser Gruppe, so dass
\weq{2}{n = \nu j}
%CONT>>
